\pagestyle{corporativo}

\begin{frontmatter}

	\title{\projecttitle}

	\author[inst1]{Jonathan Arboleda\corref{cor1}}
	\ead{proyectos2@dmlsas.com}

	\author[inst1]{Herminsul Rosero}
	\ead{herminsul.rosero@dmlsas.com}

	\cortext[cor1]{Autor correspondiente}
	\address[inst1]{DML Ingenieros Consultores S.A.S., Cali, Colombia, 760025}

	\begin{abstract}
		\thispagestyle{corporativo}
		Este informe presenta la validación técnica e hidráulica de la propuesta de profundización del trazado en el Km 16.5 del salmueroducto de BRINSA S.A., correspondiente al cruce del Humedal Arrieros mediante Perforación Horizontal Dirigida (PHD). El rediseño contempla la instalación de dos tuberías de \SI{12}{in} NPS en HDPE PE100 --- SDR 13.6 para condensados y SDR 11 para salmuera saturada --- con una longitud real de \SI{452.26}{\meter} y una profundidad vertical de \SI{15.60}{\meter}, formando una configuración de sifón invertido. El modelamiento numérico en estado estacionario desarrollado demuestra que el incremento en las pérdidas por fricción dinámicas es del \SI{1.69}{\percent} para salmuera (SDR 11) y del \SI{1.80}{\percent} para condensados (SDR 13.6), valores hidráulicamente insignificantes que no afectan la presión de succión disponible en la estación Km8. Sin embargo, se identificaron riesgos operativos críticos que requieren mitigación obligatoria: un incremento de presión estática en fondo de \SI{1.84}{bar} para la línea de salmuera que compromete el límite de diseño de la tubería SDR 13.6 (PN12.7) si la presión en superficie supera \SI{10.86}{bar}, y riesgos de sedimentación de sal y bloqueo por aire en las crestas del perfil. Adicionalmente, al quedar enterrada a \SI{15.60}{\meter} de profundidad en una zona inundable, la tubería queda sometida a esfuerzos de interacción suelo-tubería, en particular al empuje hidrostático de flotación (empuje de Arquímedes, del orden de \SI{0.77}{\kilo\newton\per\meter}) que supera el peso propio del tubo vacío, además de la pérdida de confinamiento efectivo del relleno en suelo saturado y la erosión de la cobertura en las transiciones. En consecuencia, se concluye que el tramo enterrado requiere una estrategia de diseño de ingeniería específica que garantice su integridad estructural a la profundidad de instalación, sustentada en el relleno anular estructural (lechada cemento-bentonita o CLSM) para asegurar la transferencia de carga suelo-tubo, el control de flotación mediante water ballasting durante el tiro por PHD (ASTM~F1962) y el lastre o anclaje con factor de seguridad \mbox{$FS \geq 1.25$--$1.50$} en las transiciones de cobertura reducida. Se recomienda la instalación de válvulas ventosas automáticas de triple efecto en los puntos de lanzamiento y llegada, purgas de lavado con condensado caliente y la verificación de la clase mecánica de la tubería.
	\end{abstract}

	\begin{keyword}
		Salmueroducto \sep Perforación Horizontal Dirigida \sep Simulación Hidráulica \sep Sifón Invertido \sep HDPE PE100
	\end{keyword}

\end{frontmatter}

\pagestyle{corporativo}
\thispagestyle{corporativo}
